\documentclass[a4paper,12pt]{article}
\usepackage[left=1.5cm, right=1.5cm, top=2cm,bottom=2cm]{geometry}
\usepackage{amsmath, amssymb, amsfonts}
\usepackage{multicol}
\usepackage{enumerate}
\usepackage{xcolor}

% https://tex.stackexchange.com/questions/29834/closed-square-root-symbol
\usepackage{letltxmacro}
\makeatletter
\let\oldr@@t\r@@t
\def\r@@t#1#2{%
\setbox0=\hbox{$\oldr@@t#1{#2\,}$}\dimen0=\ht0
\advance\dimen0-0.2\ht0
\setbox2=\hbox{\vrule height\ht0 depth -\dimen0}%
{\box0\lower0.4pt\box2}}
\LetLtxMacro{\oldsqrt}{\sqrt}
\renewcommand*{\sqrt}[2][\ ]{\oldsqrt[#1]{#2}}
\makeatother

\begin{document} 
  
\section*{Piape Matemática} 
\textbf{Módulo IV - Tudo é função}\\
\textbf{Exercícios Aula 01}         
\begin{multicols}{2} 

\paragraph*{1.} Para cada uma das regras abaixo, determine o domínio, contradomínio e imagem da função.
\begin{enumerate}[a)]
    \item $B = \{2,4,6,8\}$, $A = \{1,2,3,4,5\}$
    $$\begin{aligned}
        f:B&\to A\\
        x&\mapsto \frac{x}{2}
    \end{aligned}$$

    \item $A = \{1,2,3,4,5\}$, $B = \{1,2,3,\ldots, 10\}$
    $$\begin{aligned}
        f:A&\to B\\
        x&\mapsto 2x
    \end{aligned}$$
    
    \item $A = \{1, 2, 3, 4\}$, $B = \{1, 2, 3, \ldots, 20\}$
    $$\begin{aligned}
        f:A&\to B\\
        x&\mapsto x^2
    \end{aligned}$$

    \item $A = \{0, 2, \tfrac{9}{2}, 8\}$, $B = \{0,1,2,3,4,5,6\}$
    $$\begin{aligned}
        f:A&\to B\\
        x&\mapsto \sqrt{2x} + 1
    \end{aligned}$$

    \item $A = \{0, 1, 4, 5\}$,\\ $B = \{0,1,2,3,\sqrt{3},4,\sqrt{5},\sqrt{11}\}$
    $$\begin{aligned}
        f:A&\to B\\
        x&\mapsto \sqrt{2x+ 1}
    \end{aligned}$$

\end{enumerate}

\newcolumn

\paragraph*{2.} Explique o porquê as relações abaixo não representam funções.

\begin{enumerate}[a)]

    \item $A = \mathbb{N}$, $B = \mathbb{N}$
    $$\begin{aligned}
        f:A&\to B\\ 
         n&\mapsto n/2
    \end{aligned}$$

    \item $A = \{-4,-1,0,1,4\}$, $B = \{-2,-1,0,1,2\}$
    $$\begin{aligned}
        f:A&\to B\\ 
            n&\mapsto \sqrt{n}
    \end{aligned}$$
\end{enumerate}

\paragraph*{3.} Encontre o domínio das seguintes funções:
\begin{enumerate}[a)]
    \item $f(x) = 2x + \sqrt{2x - 8}$
    \item $g(x) =2^x + \dfrac{3x}{4x-6}$
    \item $h(x) = 1+\log(4x - 4)$ 
\end{enumerate}


\end{multicols}
 

\vspace*{\fill}
{\footnotesize
\paragraph*{Gabarito} \color{darkgray} \hspace*{\fill}\\
\textbf{1.} a) Domínio = B, Contradomínio = A, Imagem = $\{1,2,3,4\}$; 
b) Domínio = A, Contradomínio = B, Imagem = $\{2,4,6,8,10\}$; 
c) Domínio = A, Contradomínio = B, Imagem = $\{1,4,9,16\}$; 
d) Domínio = A, Contradomínio = B, Imagem = $\{1,3,4,5\}$;
e) Domínio = A, Contradomínio = B, Imagem = $\{1,\sqrt{3},3,\sqrt{11}\}$;\\
\textbf{2.} a) A relação não é uma função porque, por exemplo, o número 1 não tem imagem. b) A relação não é uma função porque, por exemplo, o número -4 não tem raiz quadrada.\\
\textbf{3.} a) $\text{Dom } f  = \{x\in\mathbb{R} \mid x\ge 4\}$; b)  $\text{Dom } g  = \left\{x\in\mathbb{R} \mid x \neq \tfrac{3}{2}\right\} $; c)  $\text{Dom } h  = \{x\in\mathbb{R} \mid x > 1\}$. \color{black}
}
\end{document}